\newpage

\section{Обзор литературы}
Проблема оценки потенциальной успешности предпринимательских инициатив прошла эволюцию от качественных управленческих теорий до количественных эмпирических исследований. В данном разделе рассматриваются работы, относящиеся к следующим трем направлениям: теоретические подходы, классический статистический анализ и современные количественные методы.

\subsection{Теоретические и качественные подходы}

Ранние исследователи в данной области фокусировались на качественном анализе и теории человеческого капитала. Работы данного направления, например, в рамках концепции Lean Startup~\cite{lean-startup}, рассматривают успех как абстрактную величину, зависящую от компетенций команды и соответствия продукта рынку. Исследователи выделяют такие критические факторы, как опыт основателей, размер стартового капитала и стратегию выхода на рынок. Недостатком этих работ является отсутствие строгой количественной доказательной базы. Выводы часто строятся на опросах или разборе отдельных компаний и ситуаций, что затрудняет масштабирование результатов.

\subsection{Эмпирические исследования и статистический анализ}

С накоплением исторических данных исследователи перешли к применению статистических методов. Фундаментальной работой в этой области является статья <<The evolution of new technology ventures over 20 years: Patterns of failure, merger, and survival>>~\cite{old-study}. Авторы проводят лонгитюдный анализ, отслеживая жизненный цикл технологических предприятий на протяжении двух десятилетий. Ключевым вкладом работы является разграничение понятий <<неудача>> (банкротство) и <<слияние/поглощение>>. Исследование показывает, что факторы, влияющие на выживание, динамически меняются по мере взросления компании. Однако используемые в таких работах классические статистические методы (регрессия Кокса, логистическая регрессия) часто не способны уловить сложные нелинейные зависимости, характерные для венчурного рынка.

\subsection{Современные количественные подходы}

Современный этап характеризуется наличием больших массивов данных и использованием алгоритмов для их обработки. В статье <<Predicting the Outcome of Startups: Less Failure, More Success>>~\cite{new-study} демонстрируется применение современных классификаторов для предсказания статуса стартапа. Авторы достигают высоких показателей метрик качества, значительно превосходящих традиционные статистические методы. В прикладных исследованиях сообщества аналитиков данных~\cite{kagglexai} демонстрируется использование методов визуализации важности признаков для выявления ключевых факторов успеха.

Тем не менее, для практического применения результатов инвесторами и предпринимателями зачастую достаточно наглядного эмпирического анализа, не требующего построения сложных моделей. Важно при этом учитывать методологические особенности. Так, отсутствие контроля временного фактора может приводить к ошибке выжившего --- старые компании кажутся <<успешнее>> просто потому, что у них было больше времени для достижения <<успеха>>, а те из них, кто потерпел неудачу, не попадают в рассмотрение. В данной работе используется подход визуальной аналитики и статистических тестов с контролем по когортам, позволяющий выявить и наглядно представить закономерности в данных о стартапах, избегая этой ошибки.
